\appendix

\section{The Potential of Hiding Codec Overhead with Pipeline Overlap}

In prefill--decode disaggregated serving, KV-cache compression, network transfer, and decompression can be executed as a streaming pipeline. 
This is important because the latency of a standalone codec kernel does not necessarily translate into additional end-to-end latency: if communication is the bottleneck, the encode and decode stages can be overlapped with data transfer.

Let $S$ denote the raw KV-cache size, $\rho$ the compression ratio, $G_{\mathrm{enc}}$ the compression throughput, $G_{\mathrm{dec}}$ the decompression throughput, and $B$ the physical communication bandwidth for transferring compressed bytes. 
For one pipeline chunk, the three stage times are
\[
    T_{\mathrm{enc}} = \frac{S}{G_{\mathrm{enc}}},
    \qquad
    T_{\mathrm{xfer}} = \frac{S}{\rho B},
    \qquad
    T_{\mathrm{dec}} = \frac{S}{G_{\mathrm{dec}}}.
\]
For a long stream of chunks, the steady-state pipeline time is dominated by the slowest stage:
\[
    T_{\mathrm{pipe}}
    =
    \max\left(
        \frac{S}{G_{\mathrm{enc}}},
        \frac{S}{\rho B},
        \frac{S}{G_{\mathrm{dec}}}
    \right).
\]

The codec overhead is fully hidden when the transfer stage remains the bottleneck, i.e.
    $T_{\mathrm{xfer}} \ge T_{\mathrm{enc}}
    \quad \text{and} \quad
    T_{\mathrm{xfer}} \ge T_{\mathrm{dec}}$.
Equivalently, the maximum physical communication bandwidth for fully hiding both encoding and decoding overhead is
\[
    B_{\mathrm{hide}}
    =
    \frac{\min(G_{\mathrm{enc}}, G_{\mathrm{dec}})}{\rho}.
\]
Thus, a faster codec increases the range of communication bandwidths under which SplitZip behaves as a purely bandwidth-saving transformation with no exposed codec latency.

In our implementation, SplitZip achieves a compression throughput of $G_{\mathrm{enc}} = 332.5~\mathrm{GB/s}$ and a decompression throughput of
$G_{\mathrm{dec}} = 1232.1~\mathrm{GB/s}$. Since compression is the slower codec stage, the hiding threshold is determined by $G_{\mathrm{enc}}$:
\[
    B_{\mathrm{hide}}
    =
    \frac{\min(332.5, 1232.1)}{1.32}
    =
    \frac{332.5}{1.32}
    \approx
    251.9~\mathrm{GB/s}.
\]
Therefore, SplitZip can fully hide its codec overhead as long as the physical communication bandwidth is no higher than approximately $252~\mathrm{GB/s}$.

This threshold is well above the bandwidth of common inter-node communication technologies. For example, 400~Gb/s and 800~Gb/s Ethernet correspond to approximately $50~\mathrm{GB/s}$ and $100~\mathrm{GB/s}$, respectively, and PCIe~5.0~x16 provides roughly $64~\mathrm{GB/s}$ per direction. 
Thus, under typical datacenter interconnects, SplitZip's encoding and decoding overheads can in principle be fully covered by communication time. 
In this regime, SplitZip behaves as a bandwidth-saving transformation rather than adding exposed codec latency to the end-to-end prefill--decode transfer path.

Only very high-bandwidth intra-node GPU fabrics, such as modern NVLink-class interconnects, approach or exceed this threshold. In those settings, communication may no longer be the sole bottleneck, and codec throughput can again become visible in the end-to-end critical path.